% ch02  §2.4 / §3.1–§3.3 / §4.1–§4.2（原文 963–1295 行）

\subsection{Iso-shtuka}

回忆：$F$ 是 $X$ 的函数域。

\ssszh{iso-shtuka 的张量范畴}
设 $\IsoSht(F)$ 为如下二元组 $(V, \phi)$ 构成的范畴：
\begin{itemize}
\item $V$ 是有限维 $\br F:=F\ot_{k} \ov k$-向量空间。
\item $\phi: V\isom V$ 是 $F$-线性且 $(\ov k, \Fr)$-线性的同构。也就是说，$\phi$ 是双射，且对 $v\in V$、$c\in \ov k$ 与 $a\in F$，有
$\phi(cv)=c^{q}\phi(v)$ 与 $\phi(av)=a\phi(v)$。
\end{itemize}
我们把 $\IsoSht(F)$ 中的对象称为函数域 $F$ 上的\emph{iso-shtuka}。$(V,\phi)\in \IsoSht(F)$ 的秩定义为 $\dim_{\br F}V$。于是 $\IsoSht(F)$ 在运算 $(V_{1},\phi_{1})\ot (V_{2}, \phi_{2})=(V_{1}\ot_{\br F} V_{2}, \phi_{1}\ot\phi_{2})$ 下成为刚性 $F$-线性张量范畴。$(V,\phi)$ 的对偶对象是 $(V^{\vee}=\Hom_{\br F}(V,\br F), \phi^{\vee})$，其中 $\phi^{\vee}$ 由 $\j{\phi^{\vee}(\xi),\phi(v)}=\j{\xi, v}$（$\xi\in V^{\vee}$，$v\in V$）定义。

我们应把 $F$ 上的 iso-shtuka 看作 $X$ 的通用点上的 shtuka。更确切地说，每个 $\ov k$-点 $((x_{i})_{i\in I}, \cE, \a)\in \Sht^{I}_{\GL_{n},\bK}(\ov k)$ 都对应一个如下定义的 iso-shtuka。设 $V$ 是 $\cE$ 在 $X_{\ov k}$ 的通用点 $\y=\Spec \br F$ 处的纤维。则 $\a$ 限制到 $\y$ 上给出 $\br F$-线性同构 $\a: V\isom V\ot_{\ov k,\Fr}\ov k$。对 $v\in V$ 定义 $\phi(v)=\a^{-1}(v\ot 1)$。则 $(V,\phi)$ 是 $\IsoSht(F)$ 中秩为 $n$ 的对象。

对 $F$ 上的约化群 $G$，我们把\emph{$G$-iso-shtuka} 定义为 $F$-线性张量函子
\begin{equation}\label{GisoSh}
\Rep_{F}(G)\to \IsoSht(F).
\end{equation}
若 $G$ 是分裂的，则 $\ov k$-点 $((x_{i})_{i\in I}, \cE, \a)\in \Sht^{I}_{G,\bK}(\ov k)$ 给出一个 $G$-iso-shtuka，其定义如下：它把 $G$ 的有限维 $F$-表示 $W$ 送到 iso-shtuka $(W_{\cE,\y}, \phi_{W})$，其中 $W_{\cE}=(\cE|_{U_{\ov k}})\times^{G}W$ 是由 $\cE$ 与 $W$ 相伴的 $U_{\ov k}$ 上的向量丛，$W_{\cE,\y}$ 是它的通用纤维；而 $\phi_{W}$ 是 $\a_{W}: W_{\cE}\dr {}^{\t}(W_{\cE})\cong W_{{}^{\t}\cE}$ 在 $X_{\ov k}$ 的通用点 $\y$ 处限制之逆。

\ssszh{等晶（isocrystal）}
显然，iso-shtuka 是等晶的整体类比物。设 $x\in |X|$，令 $\br F_{x}:=\wh{F_{x}^{\ur}}$ 为 $F_{x}$ 的最大非分歧扩域之完备化。设 $\Fr_{x}\in\Aut(\br F_{x})$ 是 $k_{x}$ 的自同构 $a\mapsto a^{\# k_{x}}$ 的提升。回忆：关于 $\Fr_{x}$ 的、$\br F_{x}$ 上等晶的范畴 $\IsoCrys(F_{x})$ 由如下二元组 $(V, \phi)$ 组成：
\begin{itemize}
\item $V$ 是有限维 $\br F_{x}$-向量空间。
\item $\phi: V\isom V$ 是 $(\br F_{x}, \Fr_{x})$-线性同构。也就是说，$\phi$ 是双射，且对 $v\in V$、$a\in \br F_{x}$，有
$\phi(av)=\Fr_{x}(a)\phi(v)$。
\end{itemize}
于是 $\IsoCrys(F_{x})$ 是 $F_{x}$-线性刚性张量范畴。

给定 iso-shtuka $(V,\phi)$ 与位置 $x\in|X|$，可以通过基变换作出 $\br F_{x}$ 上的等晶
\begin{equation*}
(V,\phi)\mapsto (V\ot_{\br F}\br F_{x}, \phi^{d_{x}}\ot \Fr_{x}).
\end{equation*}
这里 $d_{x}=[k_{x}:k]$。上述构造给出 $F$-线性张量函子
\begin{equation}\label{Sht fib functor}
\om_{x}: \IsoSht(F)\to \IsoCrys(F_{x}).
\end{equation}

\ssszh{单对象的分类}
回忆等晶的 Dieudonné–Manin 分类：$\IsoCrys(F_{x})$ 是半单 $F_{x}$-线性张量范畴，其单对象由斜率 $\l\in \QQ$ 分类。斜率为 $\l=m/n$（$n\ge1$，$m\in \ZZ$ 且与 $n$ 互素）的单对象 $M(\l)\in \IsoCrys(F_{x})$ 在 $\br F_{x}$ 上的维数为 $n$。对 $\l_{1}, \l_{2}\in \QQ$，$M(\l_{1})\ot M(\l_{2})$ 同构于若干个 $M(\l_{1}+\l_{2})$ 的直和。$M(\l)$ 的对偶同构于 $M(-\l)$。$M(\l)$ 的自同态代数是 $F_{x}$ 上的中心可除代数，其不变量为 $-\l\mod\ZZ\in \QQ/\ZZ$。

Drinfeld \cite[\S2]{Dr-Pet} 描述了 $\IsoSht(F)$ 作为 Abel 范畴的结构。为陈述它，设 $F^{\sep}$ 是 $F$ 的可分闭包，考虑带 $\G_{F}:=\Gal(F^{\sep}/F)$ 作用的 $\QQ$-向量空间 $F^{\sep,\times}\ot_{\ZZ}\QQ$（通常简写为 $F^{\sep,\times}\ot\QQ$）。对每个元素 $a\in F^{\sep,\times}\ot\QQ$，存在唯一的极小有限扩张 $L_{a}/F$（在 $F^{\sep}$ 内），使得 $a$ 落在 $L^{\times}_{a}\ot\QQ\to F^{\sep,\times}\ot\QQ$ 的像中。

\begin{thm}[Drinfeld \cite{Dr-Pet}]\label{th:isoSht}
\begin{enumerate}
\item 范畴 $\IsoSht(F)$ 半单。
\item $\IsoSht(F)$ 中单对象的同构类与 $F^{\sep,\times}\ot\QQ$ 上的 $\G_{F}$-轨道自然一一对应。
\item 设 $(V,\phi)$ 是 $\IsoSht(F)$ 中的单对象，其同构类对应于 $a\in F^{\sep,\times}\ot\QQ$ 的 $\G_{F}$-轨道。按上文定义 $L_{a}$，并记 $\div(a)=\sum_{y\in |L_{a}|}\ord_{y}(a)\cdot y$（对 $L_{a}$ 的位置求和）。则
\begin{enumerate}
\item $\End(V,\phi)$ 是 $L_{a}$ 上的中心可除代数，其在 $y\in |L_{a}|$ 处的不变量为 $-\ord_{y}(a)[k_{y}:k]\mod \ZZ\in \QQ/\ZZ$。
\item $\dim_{\br F}V=[L_{a}:F]d_{a}$，其中 $d_{a}\in \NN$ 是有理数集合 $\{\ord_{y}(a)[k_{y}:k]\}_{y\in |L_{a}|}$ 的最小公分母。
\item 由于 $L_{a}\subset \End(V,\phi)$，可以把 $(V,\phi)$ 看作 $\IsoSht(L_{a})$ 中的对象，于是把 $(V,\phi)$ 在位置 $y\in |L_{a}|$ 处局部化、得到 $\om_{y}(V,\phi)\in \IsoCrys(L_{a,y})$ 是有意义的。此时 $\om_{y}(V,\phi)$ 是若干斜率为 $\ord_{y}(a)$ 的单等晶的直和。
\end{enumerate}

\end{enumerate}
\end{thm}

我们给出定理第 (2) 部分一个方向的构造。由单对象 $(V,\phi)\in \IsoSht(F)$，按如下方式得到 $a\in (F^{\sep,\times}\ot\QQ)/\G_{F}$。存在有限扩张 $k_{n}/k$ 以及 $V$ 在 $F\ot_{k}k_{n}$ 上的下降 $V_{n}$，使得 $\phi$ 限制为 $\id_{F}\ot \Fr_{k_{n}/k}$-线性映射 $\phi_{n}: V_{n}\to V_{n}$。把 $a$ 定义为 $\l^{1/n}$ 的 Galois 轨道，其中 $\l\in F^{\sep,\times}$ 是 $F\ot_{k}k_{n}$-线性映射 $(\phi_{n})^{n}$ 的任一特征值。改变 $n$、下降 $V_{n}$ 与 $(\phi_{n})^{n}$ 的特征值 $\l$ 的选择，元素 $\l^{1/n}\in F^{\sep, \times}$ 只会乘上一个被 $\G_{F}$ 作用的单位根。由此得到 $F^{\sep,\times}\ot\QQ$ 中一个良定义的 $\G_{F}$-轨道。

\ssszh{张量结构}
为描述 $\IsoSht(F)$ 上的张量结构，考虑如下函数
\begin{equation*}
\chi: K_{0}(\IsoSht(F))\to \ZZ[F^{\sep, \times}\ot\QQ]^{\G_{F}}
\end{equation*}
它把对应于 $a\in F^{\sep,\times}\ot\QQ$ 的 Galois 轨道的单对象 $(V,\phi)$ 送到函数 $\chi(V,\phi): F^{\sep, \times}\ot\QQ\to \ZZ$：该函数在 $a$ 的 $\G_{F}$-轨道上支撑、在其上取常值，且总和等于 $\dim_{\br F} V$。

下面的命题描述 $\IsoSht(F)$ 中单对象如何彼此作张量积。
\begin{prop}\label{p:IsoSht tensor} 映射 $\chi$ 是单射环同态。
\end{prop}
这可以通过计算两个单对象 $V_{1},V_{2}\in \IsoSht(F)$ 的半单 $F$-代数 $\End(V_{1}\ot V_{2})\cong\End(V_{1})\ot_{ F}\End(V_{2})$ 的局部不变量来证明。

遗忘函子给出纤维函子 $\om: \IsoSht(F)\to \Vect(\br F)$，因此 $\IsoSht(F)$ 是 Tannaka 范畴。它在 $\br F$ 上的 Tannaka 群是可对角化群 $\DD_{\br F}$，特征群为 $\xch(\DD_{\br F})=F^{\sep,\times}\ot\QQ$，并带有 $\Gal(F^{\sep}/\br F)$ 的作用。

%%%%%%%%%%%%%%
%%%%%%%%%%%%%%


\section{亲戚与表亲}\label{s:sib}
本节介绍若干与 shtuka 密切相关的构造。我们可以把 Drinfeld 模看作 shtuka 的同卵双胞胎，把 Deligne–Lusztig 簇看作年幼的弟妹，而把有限域上的动机看作表亲。由于作者学识所限，我们略去了 Scholze \cite{S-Berkeley} 在过去十年发现的一个重要的混合特征表亲。

\subsection{Drinfeld 模}

我们回忆 Drinfeld 模的概念，以及 Drinfeld 模与某类特殊 shtuka 之间的词典。

\ssszh{Drinfeld 模}
设 $\infty\in |X|$，$A=\G(X-\{\infty\}, \cO_{X})$，这是 $k$ 上的 Dedekind 整环。

我们回忆秩 $n$ 的 Drinfeld $A$-模（按 \cite{Dr-Ell} 的称法，也叫椭圆模）在试验 $k$-概形 $S$ 上的概念。它是一个二元组 $(\cG, \io)$，其中 $\cG$ 是 $S$ 上的一维加性群概形（在 $S$ 上 Zariski 局部同构于 $\GG_{a,S}$），而 $\io: A\to \End_{S}(\cG)$ 是 $k$-代数同态，且满足：对每个几何点 $s=\Spec K \to S$ 以及任意 $K$-群概形同构 $\cG_{s}\cong\GG_{a,K}$，$\io$ 在 $s$-纤维上的限制
\begin{equation*}
\io_{s}: A\to \End_{K}(\cG_{s})\cong K\j{\t}
\end{equation*}
（其中 $\t\in \End_{K}(\cG_{s})\cong\End_{K}(\GG_{a,K})$ 是 Frobenius 态射 $x\mapsto x^{q}$）形如
\begin{equation*}
\io_{s}(a)=\sum_{i=0}^{-n[k_{\infty}:k]\ord_{\infty}(a)}c_{i}(a)\t^{i}, \quad \forall a\in A-\{0\}
\end{equation*}
且对 $a\ne0$，$c_{-n[k_{\infty}:k]\ord_{\infty}(a)}(a)\ne0$。

$A$ 在 $\cG$ 的李代数上的作用给出环同态 $A\to \G(S,\cO_{S})=\End_{S}(\Lie\cG)$，等价地给出映射 $\xi: S\to \Spec A=X-\{\infty\}$。这称为 Drinfeld $A$-模 $(\cG,\io)$ 的\emph{特征}。

设 $\Dr_{n}(S)$ 是 $S$ 上秩 $n$ 的 Drinfeld $A$-模构成的群胚。

\ssszh{Drinfeld 模簇}\label{sss:one leg}
为把 Drinfeld 模与 shtuka 联系起来，我们考虑例 \ref{ex:Dr} 中模问题的一个变体。

为简单起见，假设 $\infty$ 的剩余域为 $k$。设 $\Sht_{n}^{\Dr}$ 是其 $S$-点为 $(x, \cE, \{\cE(\frac{i}{n})\}_{i\in \ZZ}, \phi)$ 的模叠，其中
\begin{enumerate}
\item $x: S\to U:=X-\{\infty\}$。
\item $\cE$ 是 $X\times S$ 上秩 $n$ 的向量丛，且对任意几何点 $s\in S$，凝聚层 Euler 示性数 $\chi(X_{s}, \cE_{s})=0$（即 $\dim \upH^{0}(X_{s},\cE_{s})=\dim\upH^{1}(X_{s},\cE_{s})$）。
\item $\cE(-1)=\cE\ot\cO_{X}(-\infty)\subset \cE(-\frac{n-1}{n})\subset\cdots\subset \cE(-\frac{1}{n})\subset \cE(0)=\cE$ 是 $\cE|_{\{\infty\}\times S}$ 的一个完全旗，即对 $i=0,1,\cdots, n-1$，$\cE(-\frac{i}{n})/\cE(-\frac{i+1}{n})$ 在 $\{\infty\}\times S$ 上局部自由、秩为一。我们把 $\cE(\frac{i}{n})$ 的定义延拓到所有 $i\in \ZZ$，要求 $\cE(1+\frac{i}{n})=\cE(\frac{i}{n})\ot\cO_{X}(\infty)$。
\item $\phi: {}^{\t}\cE\to \cE(\frac{1}{n})$，其余核在 $\G(x)$ 上局部自由、秩为一；且对所有 $i=1,\cdots,n-1$，$\phi$ 把 ${}^{\t}(\cE(-\frac{i}{n}))$ 送到 $\cE(-\frac{i-1}{n})$。
\end{enumerate}

设 $I=\{1,2\}$，分解为单点集 $I_{1}=\{1\}$ 与 $I_{2}=\{2\}$。设 $\l_{1}=(0,\cdots,0,-1)$ 与 $\l_{2}=(1,0,\cdots ,0)$ 是 $\GL_{n}$ 的两个极小余权。在例 \ref{ex:Dr} 中我们定义了模叠 $\pi: \Sht^{(I_{1},I_{2}), (\l_{1},\l_{2})}_{\GL_{n}} \to X^{2}$。于是有映射
\begin{equation}\label{Dr to Sht}
\Sht_{n}^{\Dr}\to \Sht^{(I_{1},I_{2}), (\l_{1},\l_{2})}_{\GL_{n}}|_{U\times\{\infty\}}
\end{equation}
它把 $(x, \cE, \{\cE(\frac{i}{n})\}_{i\in \ZZ}, \phi)$ 送到两条leg $(x, \infty)$ 与向量丛
\begin{equation*}
\xymatrix{\cE_{0}=\cE & \cE_{1}:={}^{\t}(\cE(\frac{-1}{n}))\ar@{_{(}->}[l]_-{\phi}\ar@{^{(}->}[r] & \cE_{2}:={}^{\t}\cE}
\end{equation*}
可以证明 \eqref{Dr to Sht} 是闭嵌入；其像中的 $\ov k$-点由这样的 shtuka 组成：它在 $\infty$ 处的等晶是斜率为 $1/n$ 的超奇异等晶。

\ssszh{词典}\label{sss:dict}
据 \cite{Mum}，存在群胚的典范等价
\begin{equation*}
\Dr_{n}(S)\cong \Sht^{\Dr}_{n}(S).
\end{equation*}

我们描述两个方向的函子，更多细节见 \cite[Chapter 6]{Goss}。以下假设 $S$ 局部 Noether 环。

从 Drinfeld 模到 shtuka。设 $(\cG,\io)$ 是 $S$ 上秩 $n$ 的 Drinfeld 模。则 $M=\un\Hom_{S}(\cG,\GG_{a,S})$ 是 $S$ 上的拟凝聚层，且作为 $\cO_{S}\j{\t}=\un\End_{S}(\GG_{a,S})$-模局部自由、秩为一。此外，$M$ 通过 $\io$ 带有与 $\cO_{S}\j{\t}$-作用交换的 $A$-作用。我们将如下定义 $M$ 上一族由 $\cO_{S}$-子模构成的典范递增滤过 $M_{i}$。在 $S$ 上局部，取 $m_{0}\in M$ 作为自由 $\cO_{S}\j{\t}$-模 $M$ 的一组基。令 $M'(\frac{i}{n})$ 是 $M$ 中由 $m_{0}, \t m_{0},\cdots, \t^{i-1}m_{0}$ 生成的 $\cO_{S}$-子模。然后对滤过 $\{M'(\frac{i}{n})\}$ 作重整：定义 $M(\frac{i}{n})$ 为那些 $m\in M$，使得当 $\ord_{\infty}(a)$ 充分大（依赖于 $m$）时，$am\in M'(\frac{i}{n}+\ord_{\infty}(a))$。$m_{0}$ 的两种局部选择对 $i\gg0$ 给出相同的 $M'(\frac{i}{n})$（因为 $M$ 作为 $\cO_{S}\j{\t}$-模的任一其他基 $m_{1}$ 都形如 $m_{1}=(c_{0}+c_{1}\t+\cdots +c_{N}\t^{N})m_{0}$，其中 $c_{1}\in \cO_{S}^{\times}$，且 $c_{i}$ 对 $i\ge1$ 幂零\zy{原文此处作 $c_{1}\in\cO_{S}^{\times}$；按 $\cO_{S}\j{\t}$ 的单位群应为常数项 $c_{0}\in\cO_{S}^{\times}$，否则 $m_{1}$ 不可能是基。[依据：$\cO_{S}\j{\t}$ 中元素可逆当且仅当其常数项可逆]}）。这说明滤过 $\{M(\frac{i}{n})\}_{i\ge1}$ 与局部基 $m_{0}$ 的选择无关，从而定义了 $M$ 的一个由凝聚 $\cO_{S}$-子模构成的典范滤过。由构造可见 $\t M(\frac{i}{n})\subset M(\frac{i+1}{n})$。注意 $A$ 按在 $\infty$ 处的极点阶滤过：$A=\cup_{j\ge0}A_{j}$。令 $\Rees(A)=\op_{j\ge0}A_{j}$，则 $X=\Proj (\Rees(A))$。由构造，$A_{j}$ 把 $M(\frac{i}{n})$ 送到 $M(\frac{i}{n}+j)$。令 $N(\frac{i}{n})$ 为分次 $\Rees(A)\ot\cO_{S}$-模，其 $j$ 次分量为 $M(\frac{i}{n}+j)$。令 $\cE(\frac{i}{n})$ 为 $X\times S=\Proj_{S}(\Rees(A)\ot\cO_{S})$ 上与 $N(\frac{i}{n})$ 相伴的凝聚层。可以验证 $\cE(\frac{i}{n})$ 是 $X\times S$ 上秩 $n$ 的向量丛，且分次嵌入 $N(\frac{i}{n})\incl N(\frac{i+1}{n})$ 诱导单射 $\cE(\frac{i}{n})\incl \cE(\frac{i+1}{n})$，其余核在 $\{\infty\}\times S$ 上秩为一。最后，映射 $\phi: {}^{\t}\cE={}^{\t}\cE(0)\to \cE(\frac{1}{n})$ 来自 $\Fr_{S}$-线性映射 $\phi_{N}: N(0)\to N(\frac{1}{n})$，后者在 $j$ 次分量上是 $\t: M(j)\to M(\frac{1}{n}+j)$。

从 shtuka 到 Drinfeld 模。设 $(x, \cE, \{\cE(\frac{i}{n})\}_{i\in \ZZ}, \phi)\in \Sht^{\Dr}(S)$。令 $\pr_{S}: X\times S\to S$ 为投影。令 $M$ 为拟凝聚 $\cO_{S}$-模 $M=\varinjlim_{i}\pr_{S*}\cE(\frac{i}{n})$，它带有 $A$-作用。映射 $\phi$ 给出 $\Fr_{S}$-线性映射 $\phi_{M}: M\to M$。可以证明 $M$ 是 $\cO_{S}$ 上斜多项式环 $\cO_{S}\j{\t}$ 上的局部自由秩一模，其中 $\t$ 通过 $\phi_{M}$ 作用。由于 $\un\End_{S}(\GG_{a,S})=\cO_{S}\j{\t}$，秩一局部自由 $\cO_{S}\j{\t}$-模等同于 $\un\Aut(\GG_{a,S})\cong \cO_{S}\j{\t}^{\times}$ 的 Zariski 局部平凡主丛，也等同于 $\GG_{a}$ 在 $S$ 上的 Zariski 局部平凡形式 $\cG$。由于 $M$ 带有与 $\cO_{S}\j{\t}$-作用交换的 $A$-作用，$\cG$ 也带有 $A$-作用，这给出 Drinfeld 模的数据。

\ssszh{$t$-动机}
G. Anderson 的 $t$-动机理论 \cite{And} 是 $A=k[t]$ 情形的 Drinfeld 模的推广。一方面，若把 Drinfeld 模看作椭圆曲线的函数域类比，则 Anderson 的 $t$-动机是更一般动机的类比，包括高维 Abel 簇的类比。另一方面，\S\ref{sss:dict} 中的词典可推广到 $t$-动机：$t$-动机可以转成 $X=\PP^{1}$ 上的 shtuka，其中一条leg固定在 $\infty$，另一条leg移动。细节参见 Anderson 的原始论文 \cite{And} 与 D. Goss 的专著 \cite{Goss}。

\subsection{（仿射）Deligne–Lusztig 簇}
设 $\cB$ 是 $G$ 的旗簇，即 $G$ 的 Borel 子群的模空间。$\cB\times\cB$ 上（对角作用下的）$G$-轨道集合构成 Coxeter 群 $(W,S)$，它同构于 $G$ 的极大分裂环面所附的 Weyl 群。对 $(B,B')\in \cB^{2}$，若 $(B,B')$ 落在由 $w\in W$ 标号的 $G$-轨道中，记 $\pos(B,B')=w$；若 $(B,B')$ 落在该轨道的闭包中，记 $\pos(B,B')\le w$。

把 $G$-主丛换成 Borel 子群，可以定义 $\Hk^{(I_{1},\cdots, I_{r})}_{G}$ 的类比（至少在每个 $I_{j}$ 为单点集时）。设 $\un w=(w_{j})_{1\le j\le r}\in W^{r}$。定义 $Y^{\un w}$ 为 $\cB^{r+1}$ 中由满足 $\pos(B_{j-1},B_{j})=w_{j}$（$1\le j\le r$）的 Borel 子群 $(B_{0},\cdots, B_{r})$ 构成的子概形。类似地，定义 $Y^{\le \un w}\subset \cB^{r+1}$ 为满足 $\pos(B_{j-1},B_{j})\le w_{j}$（$1\le j\le r$）的闭子概形。

于是可以定义 $\Sht^{(I_{1},\cdots, I_{r})}_{G}$ 的类比。用如下 Cartesian 图定义 $X^{\un w}$ 与 $X^{\le \un w}$：
\begin{equation*}
\xymatrix{ X^{\un w}\ar[r]\ar[d] & Y^{\un w}\ar[d]^{(p_{0}, p_{r})} & X^{\le\un w}\ar[r]\ar[d]  & Y^{\le\un w}\ar[d]^{(p_{0}, p_{r})} \\
\cB\ar[r]^-{(\id,\Fr)} & \cB^{2} & \cB\ar[r]^-{(\id,\Fr)} & \cB^{2}
}
\end{equation*}
这里 $p_{j}$ 是第 $j$ 个投影 $\cB^{r+1}\to \cB$。注意当 $r=1$ 时，$X^{w}$ 就是 Deligne–Lusztig 簇 \cite[Definition 1.4]{DL}，它分类满足 $\pos(B,\Fr(B))=w$ 的 Borel 子群 $B$。Deligne–Lusztig 簇在 Lie 型有限群的不可约表示的构造与分类中起关键作用。

类似地，若把 $G$ 换成环路群 $LG$，$B$ 换成 Iwahori 子群 $\bI\subset LG$，$\cB$ 换成仿射旗簇 $\Fl_{G}=LG/\bI$，$W$ 换成扩张仿射 Weyl 群 $\tilW$，则可类似地定义 $\cY^{\un w}\subset \cY^{\le\un w}\subset (\Fl_{G})^{r+1}$，其中 $\un w\in (\tilW)^{r}$。再把它们与 $\Fl_{G}$ 上 Frobenius 的图相交，便得到\emph{仿射 Deligne–Lusztig 簇} $\cX^{\un w}$ 与 $\cX^{\le \un  w}$ 的特殊情形。

更一般的仿射 Deligne–Lusztig 簇由 Rapoport \cite[\S4]{Ra} 定义。为定义它们，固定 $b\in LG$，并把上一段中的 Frobenius 映射 $\Fl_{G}\to \Fl_{G}$ 换成 $b\c\Fr$，其中 $b$ 通过左平移作用。所得 $\Fl_{G}$ 的子 ind-概形记作 $\cX^{\un w}(b)$ 与 $\cX^{\le\un w}(b)$，其同构类只依赖 $b$ 的 $\Fr$-共轭类。

\subsection{有限域上的动机}
现在我们指出 shtuka 与有限域上动机之间的相似之处。主要参考文献是 Milne \cite{Milne} 与 Langlands–Rapoport \cite{LR}。

固定素数 $p$。对 $\FF_{p}$ 的每个有限扩张 $\FF_{q}$，存在 $\FF_{q}$ 上（纯）动机的 $\QQ$-线性张量范畴 $\Mot(\FF_{q})$，其 hom 空间用乘积上的圈模数值等价来定义。类似地有 $\ov\FF_{p}$ 上（纯）动机的 $\QQ$-线性张量范畴 $\Mot(\ov\FF_{p})$，它是 $\Mot(\FF_{q})$ 当 $q$ 遍历 $\FF_{p}$ 的有限扩张时的余极限。已知这些范畴半单 \cite{Jan}。

$\Mot(\ov\FF_{p})$ 作为 Abel 范畴有如下描述，它推广了有限域上 Abel 簇的 Honda–Tate 理论。回忆：权为 $n\in \ZZ$ 的 Weil $p$-数是代数数 $\pi$，它对所有素数 $\ell\ne p$ 都是 $\ell$-进单位，且 $\pi$ 的所有 阿基米德 范数都等于 $p^{n/2}$。设 $W_{p}\subset \ov \QQ$ 为各种整数权的 Weil $p$-数之集，$\mu_{\infty}\subset \ov\QQ^{\times}$ 为单位根群；后者被 Galois 群 $\G_{\QQ}=\Gal(\ov\QQ/\QQ)$ 置换。给定 $\pi\in W_{p}/\mu_{\infty}$，定义数域 $L_{\pi}=\cap_{n\in\NN}\QQ[\pi^{n}]$。

下面的定理应与定理 \ref{th:isoSht} 对照。
\begin{thm}[参见 {\cite[Theorem 2.18, 2.19]{Milne}}]\label{th:mot}
\begin{enumerate}
\item $\Mot(\ov\FF_{p})$ 中单对象的同构类与 $W_{p}/\mu_{\infty}$ 上的 $\G_{\QQ}$-轨道自然一一对应。
\item 若单动机 $X\in \Mot(\ov\FF_{p})$ 对应于权为 $w$ 的 $\pi\in W_{p}/\mu_{\infty}$ 的 Galois 轨道，则
\begin{enumerate}
\item $\End(X)$ 是数域 $L_{\pi}$ 上的中心可除代数，其在位置 $v\in |L_{\pi}|$ 处的不变量为
\begin{equation*}
\inv_{v}(\End(X))=\begin{cases}  -\ord_{v}(\pi)[k_{v}:\FF_{p}], & v\mid p \\ \frac{1}{2}, &  v\mbox{ 实，$w$ 为奇数} ; \\ 0, & \text{其他情形. } \end{cases}
\end{equation*}

\item $X$ 的秩为 $[L_{\pi}:\QQ]d_{\pi}$，其中 $d_{\pi}$ 是有理数集合 $\{\ord_{v}(\pi)[k_{v}:\FF_{p}]\}_{v|p}$ 的最小公分母。
\end{enumerate}
\end{enumerate}
\end{thm}

为描述 $\Mot(\ov\FF_{p})$ 上的张量结构，考虑如下线性函数
\begin{equation*}
\chi: K_{0}(\Mot(\ov\FF_{p}))\to \ZZ[W_{p}/\mu_{\infty}]^{\G_{\QQ}}
\end{equation*}
它把对应于 $\pi\in W_{p}/\mu_{\infty}$ 的 Galois 轨道的单动机 $X$ 送到函数 $\chi(X): W_{p}/\mu_{\infty}\to \ZZ$：该函数在 $\pi$ 的 $\G_{\QQ}$-轨道上支撑，并取常值 $d_{\pi}$。下述命题与命题 \ref{p:IsoSht tensor} 类似：

\begin{prop}[参见 {\cite[Corollary 2.20]{Milne}}] 映射 $\chi$ 是单射环同态。
\end{prop}

\ssszh{纤维函子}
对素数 $\ell\ne p$，取动机的 $\ell$-进 étale 上同调给出纤维函子
\begin{equation*}
\om_{\ell}: \Mot(\ov\FF_{p})\ot\QQ_{\ell}\to \Vect(\Ql).
\end{equation*}
对 $\ell=p$，取 $\ov\FF_{p}$ 上动机的晶体上同调给出张量函子
\begin{equation*}
\om_{p}: \Mot(\ov\FF_{p})\ot\QQ_{p}\to \IsoCrys(\QQ_{p}).
\end{equation*}
假定 Tate 猜想，$\Mot(\ov\FF_{p})\ot\RR$ 也容许到 $\RR$-向量空间的纤维函子。这些纤维函子是 \eqref{Sht fib functor} 中定义的 $\IsoSht(F)$ 上函子 $\om_{x}$ 的类比。

\ssszh{$\QQ$ 上的 iso-shtuka？}
对照定理 \ref{th:mot} 与定理 \ref{th:isoSht}，这提示应把 $\ov\FF_{p}$ 上的动机看作数域 $\QQ$ 上的 iso-shtuka，其leg在 $p$ 与 $\infty$。这一类比在 $\ov\FF_{p}$ 上椭圆曲线的动机与秩 $2$ 的 Drinfeld 模对应的 shtuka 之间尤其贴切。

在 \cite{Kot} 中，Kottwitz 把半单 $\QQ$-线性 Abel 范畴定义为 $\QQ$ 上以投射环面为带的 gerbe 的表示。猜想它对所有素数 $p$ 都把 $\Mot(\ov\FF_{p})$ 作为全子范畴包含在内。在 \cite{SchICM} 中，Scholze 建议把 Kottwitz 的范畴看作 $\IsoSht(\QQ)$（带有任意有限条leg）。

\begin{ques} 若 $\ov\FF_{p}$ 上的动机给出 $\QQ$ 上leg在 $p$ 与 $\infty$ 的 iso-shtuka，那么什么样的几何对象给出 $\QQ$ 上leg在若干素数的 iso-shtuka？
\end{ques}

%%%%%%%%%%%%%%%%%%%%%
\section{模叠的几何}\label{s:geom}
本节总结 $G$-shtuka 模叠的基本几何性质与结构。正是它们造就了这些模叠上同调上的对称性，下一节将研究这些对称性。

\subsection{可表示性}
关于 shtuka 模叠的几何，最基本的结果如下。
\begin{thm}[Varshavsky {\cite[Proposition 2.16(a)]{Var}}] 设 $I=I_{1}\coprod\cdots\coprod I_{r}$ 是有限集的分解，且 $\un\l=(\l_{i})_{i\in I}\in (\xcoch(T)^{+})^{I}$。则叠 $\Sht^{(I_{1},\cdots, I_{r}), \le \un\l}_{G,\bK}$ 是 $k$ 上局部有限型的 Deligne–Mumford 叠。
\end{thm}

众所周知，$\Bun_{G,\bK}$ 是 $k$ 上局部有限型的光滑 Artin 叠。此外，映射 $p_{i}: \Hk_{G,\bK}^{(I_{1},\cdots, I_{r}), \le \un\l}\to \Bun_{G,\bK}$ 有限型。由定义 $\Sht^{(I_{1},\cdots, I_{r}), \le\un\l}_{G,\bK}$ 的图 \eqref{defn Sht} 可见，它是 $k$ 上局部有限型的 Artin 叠。

\ssszh{Harder–Narasimhan 截断}\label{sss:HN}
可以如下用 $k$ 上有限型的开子叠覆盖 $\Sht^{(I_{1},\cdots, I_{r}), \le \un\l}_{G,\bK}$。

先考虑无水平结构的情形。设 $\xcoch(T)^{+}_{\QQ}$ 是 $\xcoch(T)_{\QQ}$ 中的支配锥。设 $\HN\subset \xcoch(T)^{+}_{\QQ}$ 为满足如下整性条件的 $\mu\in \xcoch(T)_{\QQ}^{+}$ 之集：若 $L_{\mu}$ 是以 $\{\a|\j{\a,\mu}=0\}$ 为根的标准 Levi 子群，则 $\mu\in \xcoch(L_{\mu,\ab})_{\QQ}=\xcoch(ZL_{\mu})_{\QQ}\subset\xcoch(T)_{\QQ}$（这里 $L_{\mu,\ab}$ 是 $L_{\mu}$ 的交换化）。每个几何点 $\cE\in \Bun_{G}$ 都有关于标准抛物子群 $P\subset G$ 的典范 $P$-约化 $\cE_{P}$（推广向量丛的 Harder–Narasimhan 滤过），使得：设 $L$ 是 $P$ 的 Levi 商，则诱导的 $L_{\ab}$-主丛 $\cE_{L_{\ab}}=\cE_{P}/(U_{P}[L,L])$ 的次数 $\deg(\cE_{L_{\ab}})\in \xcoch(L_{\ab})$ 满足：对 $U_{P}$ 中所有根 $\a$，$\j{\a, \deg(\cE_{L_{\ab}})}>0$。因此每个几何点 $\cE\in \Bun_{G}$ 都有良定义的 Harder–Narasimhan 点 $\deg(\cE_{L_{\ab}})\in \HN$。

考虑 $\HN$ 上的偏序：$\mu\le \mu'$ 当且仅当 $\mu'- \mu $ 是单余根的 $\QQ_{\ge0}$-线性组合。存在开子叠 ${}^{\le\mu}\Bun_{G}\subset \Bun_{G}$，其几何点的 HN 点 $\le \mu$。已知 ${}^{\le\mu}\Bun_{G}$ 在 $k$ 上有限型。

容易看出，记录 $\cE_{0}$ 的遗忘映射 $p_{0}: \Sht^{(I_{1},\cdots, I_{r}), \le \un\l}_{G}\to \Bun_{G}$ 有限型。现对 $\mu\in \HN$，令 ${}^{\le \mu}\Sht^{(I_{1},\cdots, I_{r}), \le \un\l}_{G}=p_{0}^{-1}({}^{\le \mu}\Bun_{G})$，则 ${}^{\le \mu}\Sht^{(I_{1},\cdots, I_{r}), \le \un\l}_{G}$ 在 ${}^{\le \mu}\Bun_{G}$ 上有限型，从而在 $k$ 上有限型。当 $\mu\in \HN$ 变动时，$\{{}^{\le \mu}\Sht^{(I_{1},\cdots, I_{r}),\le\un\l}_{G}\}_{\mu\in\HN}$ 覆盖 $\Sht^{(I_{1},\cdots, I_{r}),\le\un\l}_{G}$。

对一般的水平结构 $\bK=(\bK_{x})_{x\in \Sig}$，可取 $n_{x}\in\NN$，使得 $\bK_{x}$ 包含主同余子群 $L^{(n_{x})}_{x}G$。令 $D=\sum_{x\in \Sig}n_{x}x$ 与 $\bK_{D}=(L^{(n_{x})}_{x}G)_{x\in \Sig}$，则 $\Sht^{(I_{1},\cdots, I_{r})}_{G,\bK_{D}}\to \Sht^{(I_{1},\cdots, I_{r})}_{G,\bK}$ 是有限 étale 映射（见 \S\ref{ss:change level}），因此只需用 $k$ 上有限型的开子叠覆盖 $\Sht^{(I_{1},\cdots, I_{r})}_{G,\bK_{D}}$。有遗忘映射 $\Sht^{(I_{1},\cdots, I_{r}),\le\un\l}_{G,\bK_{D}}\xr{p_{0}}\Bun_{G,D}\to \Bun_{G}$。对 $\mu\in \HN$，取 ${}^{\le \mu}\Sht^{(I_{1},\cdots, I_{r}),\le\un\l}_{G,\bK_{D}}$ 为 ${}^{\le\mu}\Bun_{G}$ 的原像，则 ${}^{\le \mu}\Sht^{(I_{1},\cdots, I_{r}),\le\un\l}_{G,\bK_{D}}$ 在 $\Sht^{(I_{1},\cdots, I_{r}),\le\un\l}_{G,\bK_{D}}$ 中是开的，且在 $k$ 上有限型。当 $\mu\in \HN$ 变动时，${}^{\le \mu}\Sht^{(I_{1},\cdots, I_{r}),\le\un\l}_{G,\bK_{D}}$ 覆盖 $\Sht^{(I_{1},\cdots, I_{r}),\le\un\l}_{G,\bK_{D}}$。

\medskip\noindent Varshavsky \cite[Prop.2.16(1)]{Var} 实际上证明了更精确的结果：每个有限型开子叠 ${}^{\le\mu}\Sht^{(I_{1},\cdots, I_{r}), \le\un\l}_{G,\bK}$ 都是一个拟射影概形被一个有限（常）群作用后的商。论证思路是：通过过渡到更深的水平结构 $\bK'_{x}$（仍在 $x\in \Sig$ 处），其在 $\bK_{x}$ 中正规，可使 ${}^{\le\mu}\Bun_{G,\bK'}$ 成为拟射影概形。于是由图 \eqref{defn Sht} 清楚可见，${}^{\le\mu}\Sht^{(I_{1},\cdots, I_{r}), \le\un\l}_{G,\bK'}$ 是拟射影概形。由 \S\ref{ss:change level} 的讨论，映射 ${}^{\le\mu}\Sht^{(I_{1},\cdots, I_{r}), \le\un\l}_{G,\bK'}\to {}^{\le\mu}\Sht^{(I_{1},\cdots, I_{r}), \le\un\l}_{G,\bK}$ 是有限群 $\prod_{s\in \Sig}K_{x}/K'_{x}$ 的主齐性空间，其中 $K_{x}=\bK_{x}(k_{x})$，$K'_{x}=\bK'_{x}(k_{x})$。

\subsection{局部模型}
人们希望了解 shtuka 模叠更精细的几何信息，例如其光滑性与维数。这些可以借助仿射 Grassmann 簇中的仿射 Schubert 簇给出的局部模型来理解。

设 $I=I_{1}\coprod\cdots\coprod I_{r}$ 是分成\emph{单点集}的分解。回忆相对位置映射 \eqref{ev disj}：
\begin{equation}\label{ev Hk singleton}
\ev^{(I_{1},\cdots, I_{r})}_{\Hk}: \Hk^{(I_{1},\cdots, I_{r})}_{G,\bK}\to (\QGr/\Aut(D))^{I}.
\end{equation}
注意由于每个 $I_{j}$ 都是单点集，该映射在整个 $U^{I}$ 上有定义。考虑复合
\begin{equation}\label{ev Sht singleton}
\ev^{(I_{1},\cdots, I_{r})}_{\Sht}: \Sht^{(I_{1},\cdots, I_{r})}_{G,\bK}\xr{} \Hk^{(I_{1},\cdots, I_{r})}_{G,\bK}\xr{\ev^{(I_{1},\cdots, I_{r})}_{\Hk}} (\QGr/\Aut(D))^{I}.
\end{equation}

\begin{prop}\label{p:fm sm} 设 $I=I_{1}\coprod\cdots\coprod I_{r}$ 是分成单点集的分解。则映射 \eqref{ev Hk singleton} 与 \eqref{ev Sht singleton} 形式光滑。
\end{prop}
\begin{proof}[证明梗概]

先证 \eqref{ev Sht singleton} 形式光滑。

回忆 $D=\Spec k\tl{t}$。设 $\wh U\to U$ 是典范 $\Aut(D)$-主丛 $(x,t)$，其中 $t$ 是 $x\in U$ 处的局部参数。则我们有映射
\begin{equation*}
\wh\ev: \wh\Sht^{(I_{1},\cdots, I_{r})}_{G,\bK}:=\wh U^{I}\times_{U^{I}}\Sht^{(I_{1},\cdots, I_{r})}_{G,\bK}\to  \wh U^{I}\times \QGr^{I}
\end{equation*}
使得 $\ev^{(I_{1},\cdots, I_{r})}$ 由投影到 $\QGr^{I}$ 并在两边对 $\Aut(D)^{I}$ 取商得到。因此只需证明 $\wh\ev$ 形式光滑。

为记号简洁，只处理无水平结构的情形，一般情形类似。可把 $I$ 与 $\{1,\cdots, r\}$ 等同，使 $I_{j}=\{j\}$。设 $R'$ 是 $k$-代数，$J\subset R'$ 是平方为零的理想，$R=R'/J$。假设给定 $R'$-点 $(x'_{j}, t'_{j})\in \wh X(R')$（$1\le j\le r$）、$R'$-点 $\z'_{j}=(\cF'_{j}\dxr{} \cG'_{j})\in \QGr(R')$（$1\le j\le r$，其中 $(\cF'_{j}, \cG'_{j})$ 是 $D_{R'}$ 上的一对 $G$-主丛，并在 $D^{\times}_{R'}$ 上带同构），以及 $\wh \Sht^{(I_{1},\cdots, I_{r})}_{G}$ 的 $R$-点 $\xi=((x_{j},t_{j})_{1\le j\le r}, (\cE_{j})_{0\le j\le r}, (\a_{j})_{1\le j\le r}, \io)$，使得 $\wh\ev(\xi)=((x'_{j}, t'_{j},\z'_{j})|_{\Spec R})_{1\le j\le r}$。特别地，$(x_{j},t_{j})=(x'_{j}, t'_{j})|_{\Spec R}$。我们希望把 $\xi$ 提升为 $\wh \Sht^{(I_{1},\cdots, I_{r})}_{G}$ 的 $R'$-点 $\xi'$，使得 $\wh\ev(\xi')=(x'_{j}, t'_{j}, \z'_{j})_{1\le j\le r}$。这意味着需要把 $\cE_{j}$ 提升为 $X_{R'}$ 上的 $G$-主丛 $\cE'_{j}$（$0\le j\le r$），把同构 $\a_{j}$ 提升为 $\a'_{j}: \cE'_{j-1}|_{X_{R'}-\G(x'_{j})}\isom \cE'_{j}|_{X_{R'}-\G(x'_{j})}$（$1\le j\le r$），并把同构 $\io: \cE_{r}\cong(\id_{X}\times \Fr_{R})^{*}\cE_{0}$ 提升为同构 $\io': \cE'_{r}\cong(\id_{X}\times \Fr_{R'})^{*}\cE'_{0}$。

首先，$\cE'_{r}$ 已由 $\cE_{0}$ 如下确定。由于 $J^{2}=0$ 是平方零加厚，$\Fr_{R'}$ 分解为 $R'\surj R\xr{\ov\Fr_{R'}} R'$。设 $\ph: \Spec R'\to \Spec R$ 是由 $\ov\Fr_{R'}$ 给出的映射。若 $\cE'_{0}$ 是 $\cE_{0}$ 到 $X_{R'}$ 的任一扩张，则 $(\id_{X}\times \Fr_{R'})^{*}\cE'_{0}=(\id_{X}\times \ph)^{*}\cE_{0}$。因此，为得到所需同构 $\io': \cE'_{r}\cong {}^{\t}\cE'_{0}$，必须把 $\cE'_{r}$ 定义为 $(\id_{X}\times \ph)^{*}\cE_{0}$。

接着把沿 $\G(x_{j})$ 的修正 $\a_{j}: \cE_{j-1}\dxr\cE_{j}$ 延拓为沿 $\G(x'_{j})$ 的修正 $\a'_{j}: \cE'_{j-1}\dxr\cE'_{j}$。我们从 $j=r$ 开始归纳进行。每一步给定 $X_{R'}$ 上的 $\cE'_{j}$，试图构造 $\cE'_{j-1}$ 与 $\a'_{j}$。用 $t'_{j}$ 把 $D_{x'_{j}}$ 与 $D_{R'}$ 等同，从而也把 $D_{x_{j}}$ 与 $D_{R}$ 等同。同构 $\wh\ev(\xi)=(x_{j}, t_{j},\z'_{j}|_{\Spec R})_{1\le j\le r}$ 给出图之间的同构
\begin{equation*}
\xymatrix{\cE_{j-1}|_{D_{x_{j}}}\ar[d]^{\e_{j-1}}_{\wr}\ar@{-->}[r]^{\a_{j}} & \cE_{j}|_{D_{x_{j}}}\ar[d]^{\e_{j}}_{\wr}\\
\cF'_{j}|_{D_{R}}\ar@{-->}[r]^{\z_{j}} & \cG'_{j}|_{D_{R}}
}
\end{equation*}
先把同构 $\e_{j}$ 提升为 $\e'_{j}: \cE'_{j}|_{D_{x'_{j}}}\cong \cG'_{j}$（之所以可行，是因为 $L^{+}G$ 形式光滑）。然后把 $\cE'_{j-1}$ 定义为 $\cE'_{j}|_{X_{R'}-\G(x'_{j})}$ 与 $\cF'_{j}$ 沿同构 $\cE'_{j}|_{D^{\times}_{x'_{j}}}\xr{\e'_{j}} \cG'_{j}|_{D^{\times}_{R'}}\xr{\z_{j}^{-1}}\cF'_{j}|_{D^{\times}_{R'}}$ 的粘合。于是有典范同构 $\cE'_{j-1}|_{X_{R}}\cong \cE_{j-1}$，因为后者由 $\cE_{j}$ 与 $\cF'_{j}|_{D_{R}}$ 沿 $\cE_{j}|_{D^{\times}_{x_{j}}}\xr{\e_{j}}\cG'_{j}|_{D^{\times}_{R}}\xr{\z_{j}^{-1}}\cF'_{j}|_{D^{\times}_{R}}$ 粘合而成。这就完成了 $\cE'_{j}$ 的归纳构造。因此 \eqref{ev Sht singleton} 形式光滑。

为证 \eqref{ev Hk singleton} 形式光滑，论证类似：只需从 $\cE_{r}$ 到 $X_{R'}$ 的任一提升 $\cE'_{r}$ 出发，其存在性来自 $\Bun_{G}$ 的光滑性。
\end{proof}

过渡到有界版本，命题 \ref{p:fm sm} 中证明的形式光滑性蕴含有界版本相对位置映射的光滑性。更强地，仿射 Grassmann 簇在如下意义下是模叠的 étale 局部模型。

\begin{prop}[{\cite[Theorem 2.20]{Var}}，{\cite[Theorem 2.11]{VLaff}}]\label{p:local model} 设 $I=I_{1}\coprod\cdots\coprod I_{r}$ 是有限集 $I$ 的分解，$\un\l\in (\xcoch(T)^{+})^{I}$。

\begin{enumerate}
\item 叠 $\Sht^{(I_{1},\cdots, I_{r}), \le \un\l}_{G,\bK}$ 在 $U^{I}$ 上 étale 局部同构于 $\Gr^{(I_{1},\cdots, I_{r}), \le\un\l}_{G}$，后者又 étale 局部同构于乘积 $\prod_{j=1}^{r}\Gr^{I_{j},\le\un\l_{j}}_{G}$（其中 $\un\l_{j}=\{\l_{i}\}_{i\in I_{j}}$）。
\item 特别地，若 $\un\l$ 是 $G$-容许的，则
\begin{equation*}
\dim \Sht^{(I_{1},\cdots, I_{r}),\le\un\l}_{G,\bK}=\sum_{i\in I}(\j{2\r, \l_{i}}+1).
\end{equation*}
\item 设对每个 $j$，$I_{j}=I_{j,1}\coprod\cdots\coprod I_{j,s_{j}}$ 是分解，并令 $I=I'_{1}\coprod\cdots\coprod I'_{r'}$ 是 $I=I_{1}\coprod\cdots\coprod I_{r}$ 由此得到的加细。则态射 $\pi^{(I'_{1},\cdots, I'_{r'}),\le\un\l}_{(I_{1},\cdots, I_{r})}: \Sht^{(I'_{1},\cdots, I'_{r'}), \le \un\l}_{G,\bK}\to \Sht^{(I_{1},\cdots, I_{r}), \le \un\l}_{G,\bK}$ 在目标上 étale 局部同构于自然映射之积 $\prod_{j=1}^{r}\Gr^{(I_{j,1},\cdots, I_{j,s_{j}}),\le\un\l_{j}}_{G}\to \prod_{j=1}^{r}\Gr^{I_{j},\le\un\l_{j}}_{G}$。
\end{enumerate}
\end{prop}
证明最终归结为一个切向计算，并用到了 Frobenius 映射的微分为零。

若 $\l_{i}$ 都是极小的，则 $\Gr^{\le\l_{i}}_{G}$ 光滑，从而 $\Gr^{(I_{1},\cdots, I_{r}), \le\un\l}_{G}$ 在 $U^{I}$ 上光滑。此时命题 \ref{p:local model} 给出：
\begin{cor}\label{c:Sht sm} 设 $I=I_{1}\coprod\cdots\coprod I_{r}$ 是分成单点集的分解，且对所有 $i\in I$，$\l_{i}\in \xcoch(T)^{+}$ 是极小支配余权。则 $\Sht^{(I_{1},\cdots, I_{r}),\le\un\l}_{G,\bK}$ 在 $U^{I}$ 上光滑，纯相对维数为 $\sum_{i\in I}\j{2\r, \l_{i}}$。
\end{cor}
